%! Unit 1: Asking Questions with Data: Study Design — Summative Assessment — Unit Test (KEY)
\documentclass[10pt]{article}
\usepackage{saar-article}
\usepackage{saar-key}   % pulls in -boxes; do NOT also load -boxes

% Work blocks on a test need handwriting room, not typeset spacing. This moves
% the blank and the key together, so it cannot open a page mismatch.
\setlength{\workrowsep}{5pt}

% Part-divider strip
\newcommand{\parthead}[1]{%
  \vspace{6pt}\noindent
  \begin{headlinebox}{cerulean}{\color{white}\bfseries #1}\end{headlinebox}%
  \vspace{4pt}}

% NO teachernote here: this key must stay the same length as its blank, and the
% conventions put a test's rationale where students cannot reach it. Answer
% rationale and Part D scoring live on page 2 of unit01/unit_cover_key/main.tex.

\begin{document}

\pageheader{Unit 1: Asking Questions with Data: Study Design}{Unit Test --- Answer Key}
\namedateperiod

\vspace{0.06in}
\noindent\textbf{Instructions:} Show all work. No notes unless your teacher says so.
\textbf{100 points:} Part A $14$, Part B $16$, Part C $40$, Part D $30$.

% ======================================================================
\parthead{Part A --- Vocabulary \hfill 14 questions, 1 point each}

\small
\textbf{Set 1 --- Samples and Bias.} Write the letter of the matching description
next to each term.

\vspace{2pt}
\begin{minipage}[t]{0.36\linewidth}
\begin{enumerate}[leftmargin=1.9em, itemsep=3.5pt, topsep=1pt]
  \item \ans{A} Census
  \item \ans{F} Parameter
  \item \ans{B} Statistic
  \item \ans{C} Convenience sample
  \item \ans{G} Systematic sample
  \item \ans{D} Simple random sample
  \item \ans{E} Nonresponse
\end{enumerate}
\end{minipage}\hfill
\begin{minipage}[t]{0.61\linewidth}
\begin{enumerate}[label=\Alph*., leftmargin=1.9em, itemsep=3.5pt, topsep=1pt]
  \item Every individual in the population is measured.
  \item A number that describes a sample.
  \item Individuals are chosen because they are the easiest ones to reach.
  \item Every group of the chosen size is equally likely to be the one selected.
  \item People who were chosen for the sample never answered.
  \item A number that describes a population.
  \item Start at a random individual on the list, then take every $k$th one
        after that.
\end{enumerate}
\end{minipage}

\vspace{6pt}
\textbf{Set 2 --- Studies and Experiments.} Same directions.

\vspace{2pt}
\begin{minipage}[t]{0.36\linewidth}
\begin{enumerate}[leftmargin=1.9em, itemsep=3.5pt, topsep=1pt, start=8]
  \item \ans{C} Observational study
  \item \ans{E} Experiment
  \item \ans{F} Control group
  \item \ans{A} Randomization
  \item \ans{D} Replication
  \item \ans{G} Confounding variable
  \item \ans{B} Placebo effect
\end{enumerate}
\end{minipage}\hfill
\begin{minipage}[t]{0.61\linewidth}
\begin{enumerate}[label=\Alph*., leftmargin=1.9em, itemsep=3.5pt, topsep=1pt]
  \item Using chance to decide which group each subject lands in.
  \item Subjects respond to being chosen and treated, not to the treatment itself.
  \item The researcher measures or surveys without assigning anyone to a group.
  \item Enough units in each group that a result is not one person's luck.
  \item The researcher assigns individuals to groups and imposes a condition.
  \item The baseline group, which gets no treatment or the usual condition.
  \item A second difference between the groups, so you cannot tell which one
        produced the result.
\end{enumerate}
\end{minipage}
\normalsize

% ======================================================================
\boxguard[8]
\parthead{Part B --- Multiple Choice \hfill 8 questions, 2 points each}

\small
\begin{enumerate}[leftmargin=2.2em, itemsep=5pt, topsep=2pt, start=15]

\item Oakmont Middle School records each player's \textbf{jersey number} ($7$,
$12$, $23$, \ldots). This variable is
\begin{enumerate}[label=(\Alph*), leftmargin=1.8em, itemsep=1pt, topsep=1pt]
  \item quantitative, because the values are numbers.
  \item quantitative, because you can compute a mean of them.
  \item categorical, because the numbers are labels for players.
        \textcolor{keyred}{\textbf{$\leftarrow$ correct}}
  \item neither, because jersey numbers are assigned at random.
\end{enumerate}

\item In a random sample of $80$ Northgate members, $35\%$ said they drive to the
center. That $35\%$ is a
\begin{enumerate}[label=(\Alph*), leftmargin=1.8em, itemsep=1pt, topsep=1pt]
  \item statistic, because it describes the $80$ members who were surveyed.
        \textcolor{keyred}{\textbf{$\leftarrow$ correct}}
  \item statistic, because the sample was chosen at random.
  \item parameter, because it is written as a percent.
  \item parameter, because it will be used to describe all the members.
\end{enumerate}

\item Which sentence describes a \textbf{cluster} sample?
\begin{enumerate}[label=(\Alph*), leftmargin=1.8em, itemsep=1pt, topsep=1pt]
  \item Some individuals are chosen at random from every group.
  \item Every $k$th individual on the list is chosen after a random start.
  \item Whoever happens to be nearby is asked.
  \item A few whole groups are chosen at random, and everyone in them is studied.
        \textcolor{keyred}{\textbf{$\leftarrow$ correct}}
\end{enumerate}

\item A survey is mailed to all $800$ households in a neighborhood, and only
$200$ send it back. The $600$ who threw it away disagree with the ones who
answered. This is an example of
\begin{enumerate}[label=(\Alph*), leftmargin=1.8em, itemsep=1pt, topsep=1pt]
  \item undercoverage, because those households were never on the list.
  \item nonresponse, because those households were chosen but did not answer.
        \textcolor{keyred}{\textbf{$\leftarrow$ correct}}
  \item sampling variability, because samples differ from each other.
  \item a response bias, because the question was worded badly.
\end{enumerate}

\item A biased survey is repeated with the same method and \textbf{twice as many
responses}. The new percent will most likely
\begin{enumerate}[label=(\Alph*), leftmargin=1.8em, itemsep=1pt, topsep=1pt]
  \item move halfway toward the true value.
  \item move all the way to the true value.
  \item stay about the same, because the method did not change.
        \textcolor{keyred}{\textbf{$\leftarrow$ correct}}
  \item become impossible to predict in any direction.
\end{enumerate}

\item In a random sample of Oakmont students, members of the Homework Club pass
\emph{fewer} classes than non-members. This result shows
\begin{enumerate}[label=(\Alph*), leftmargin=1.8em, itemsep=1pt, topsep=1pt]
  \item an association; a lurking variable could explain the whole thing.
        \textcolor{keyred}{\textbf{$\leftarrow$ correct}}
  \item that the Homework Club causes students to pass fewer classes.
  \item that passing classes causes students to avoid the Homework Club.
  \item nothing at all, because the sample was chosen at random.
\end{enumerate}

\item In an experiment, the purpose of \textbf{random assignment} is to
\begin{enumerate}[label=(\Alph*), leftmargin=1.8em, itemsep=1pt, topsep=1pt]
  \item make the two groups exactly the same size.
  \item let the subjects pick the group they would rather be in.
  \item guarantee that the treatment works.
  \item make the two groups alike at the start, so the treatment is the only
        difference left. \textcolor{keyred}{\textbf{$\leftarrow$ correct}}
\end{enumerate}

\item A researcher wants to know whether growing up in a home with heavy mold
causes asthma. No experiment is run. The best reason is that the experiment
would be
\begin{enumerate}[label=(\Alph*), leftmargin=1.8em, itemsep=1pt, topsep=1pt]
  \item impractical, because volunteers are hard to find.
  \item unethical, because it would assign children to live with mold.
        \textcolor{keyred}{\textbf{$\leftarrow$ correct}}
  \item impossible, because asthma cannot be measured.
  \item unnecessary, because observational studies prove causation.
\end{enumerate}

\end{enumerate}
\normalsize

% ======================================================================
\boxguard[8]
\parthead{Part C --- Short Answer \& Computation \hfill 8 questions, 5 points each}

\small
\begin{enumerate}[leftmargin=2.2em, itemsep=9pt, topsep=2pt, start=23]

% ---------------------------------------------------------------- C1 (1.1)
\item \textbf{Oakmont Middle School} has $600$ students. For each student the
principal records: grade level, minutes spent reading last night, whether the
student walks to school, and favorite elective.

\vspace{2pt}
(a) Who are the \textbf{individuals} in this data set? \ans{The 600 students}

\vspace{3pt}
(b) Classify each variable.

\vspace{2pt}
\renewcommand{\arraystretch}{1.25}
\begin{tabularx}{\linewidth}{@{} X c @{}}
\rowcolor{frost}
\textbf{Variable} & \textbf{Categorical or quantitative?} \\
\midrule
Grade level & \ans{Categorical} \\
Minutes of reading last night & \ans{Quantitative} \\
Walks to school (yes / no) & \ans{Categorical} \\
Favorite elective & \ans{Categorical} \\
\end{tabularx}

\vspace{4pt}
(c) The principal's first draft question is \emph{``Do students read?''}
Explain why that is not a statistical question, then rewrite it so that it is.

\ansline{It names no measurement and every student answers it the same way ---
yes or no.}
\ansline{``How many minutes did Oakmont students spend reading last night?''}

% ---------------------------------------------------------------- C2 (1.2)
\boxguard[16]
\item \textbf{Northgate Recreation Center} has $1{,}200$ members. The director
surveys $80$ members chosen at random; $28$ of them say they drive to the center.

\vspace{2pt}
(a) Population: \ans{All 1,200 members} \quad Sample size: \ans{$80$}

\vspace{3pt}
Sample: \ans{The 80 members who were surveyed}

\vspace{3pt}
(b) What percent of the sample drives? About how many of all $1{,}200$ members
does that estimate?

\begin{work}
  \dfrac{28}{80} &= 0.35 = 35\% \\
  0.35 \times 1200 &= 420 \text{ members}
\end{work}

(c) Is the $35\%$ a parameter or a statistic? Say how you know.

\ansline{A statistic --- it describes the 80-member sample, not all 1,200 members.}

(d) Give one reason a \textbf{census} of all $1{,}200$ members is impractical here.

\ansline{Cost --- surveying every one of 1,200 members would take far too much time
and money.}

% ---------------------------------------------------------------- C3 (1.3)
\boxguard[18]
\item \textbf{Westfield High School} has $800$ students and wants a
\textbf{stratified} random sample of $50$, allocated in proportion to grade size.

\vspace{2pt}
(a) What fraction of the school is being sampled?

\begin{work}
  \dfrac{50}{800} &= \dfrac{1}{16}
\end{work}

(b) Show the work for the \textbf{11th grade} row, then complete the table.

\begin{work}
  176 \times \dfrac{1}{16} &= 11 \text{ students}
\end{work}

\vspace{-2pt}
\renewcommand{\arraystretch}{1.2}
\begin{tabularx}{\linewidth}{@{} X c c c c | c @{}}
\rowcolor{frost}
\textbf{Grade} & 9 & 10 & 11 & 12 & \textbf{Total} \\
\midrule
Students in grade & $256$ & $224$ & $176$ & $144$ & $800$ \\
Number sampled & \ans{$16$} & \ans{$14$} & \ans{$11$} & \ans{$9$} & \ans{$50$} \\
\end{tabularx}

% ---------------------------------------------------------------- C4 (1.3)
\boxguard[18]
\item \textbf{Pinehurst Apartments} has $800$ households on a numbered list.
The manager wants a \textbf{systematic} random sample of $40$.

\vspace{2pt}
(a) Find the interval $k$.

\begin{work}
  k &= 800 \div 40 \\
  k &= 20
\end{work}

(b) The random start is household $6$. List the first four households selected:

\vspace{2pt}
\ans{$6$} \quad \ans{$26$} \quad \ans{$46$} \quad \ans{$66$}

\vspace{4pt}
(c) Name the sampling technique each manager used.

\vspace{2pt}
\renewcommand{\arraystretch}{1.25}
\begin{tabularx}{\linewidth}{@{} X c @{}}
\rowcolor{frost}
\textbf{What the manager did} & \textbf{Technique} \\
\midrule
Split the list by unit type, then drew households at random from
every type. & \ans{Stratified} \\
Asked the households on the two floors he manages himself. & \ans{Convenience} \\
Drew $5$ of the $40$ buildings at random and surveyed everyone in them. & \ans{Cluster} \\
\end{tabularx}

% ---------------------------------------------------------------- C5 (1.4)
\boxguard[20]
\item \textbf{Lakeside Farmers Market} has $800$ Saturday shoppers and wants to
know what fraction would come to a Wednesday market. Two studies were run.

\vspace{1pt}
\textbf{Method 1 --- clipboard at the market's own booth.} $120$ people signed
it; $96$ said yes. \\
\textbf{Method 2 --- random sample of shoppers.} $150$ were asked; $45$ said yes.

\vspace{2pt}
(a) Find the percent saying yes for each method.

\begin{work}
  \dfrac{96}{120} &= 0.80 = 80\% \\
  \dfrac{45}{150} &= 0.30 = 30\%
\end{work}

(b) Estimate the number of all $800$ Saturday shoppers each method points to.

\begin{work}
  0.80 \times 800 &= 640 \text{ shoppers} \\
  0.30 \times 800 &= 240 \text{ shoppers}
\end{work}

(c) One method is biased. Name it, name the \textbf{type} of bias, and say which
\textbf{direction} it pushes the estimate.

\ansline{Method 1. Undercoverage --- only the market's most loyal shoppers were at
that booth,}
\ansline{so it overestimates: $640$ shoppers against $240$.}

% ---------------------------------------------------------------- C6 (1.5)
\boxguard[20]
\item A random sample of $200$ Oakmont students is taken. Of the $80$ who belong
to the \textbf{Homework Club}, $44$ passed all their classes. Of the $120$ who do
not belong, $84$ passed all their classes. Nobody was assigned --- students joined.

\vspace{2pt}
(a) Find each group's pass percent, and the gap between them.

\begin{work}
  \dfrac{44}{80} &= 0.55 = 55\% \\
  \dfrac{84}{120} &= 0.70 = 70\% \\
  70\% - 55\% &= 15 \text{ percentage points}
\end{work}

(b) Explanatory variable: \ans{Homework Club membership} \quad Response: \ans{Passed all classes}

\vspace{3pt}
(c) Using all $200$ students, find the overall pass percent and estimate how
many of all $600$ students that is.

\begin{work}
  \dfrac{44 + 84}{200} &= \dfrac{128}{200} = 0.64 = 64\% \\
  0.64 \times 600 &= 384 \text{ students}
\end{work}

(d) Name one \textbf{lurking variable} that could produce this result, and say
which way it would push the numbers.

\ansline{Students who were already struggling are the ones who join the Homework
Club.}
\ansline{That pushes the members' pass rate down without the club causing it.}

% ---------------------------------------------------------------- C7 (1.6)
\boxguard[20]
\item Oakmont now runs an \textbf{experiment}. Of $100$ volunteers, a random
number generator sends $50$ to required Homework Club meetings and $50$ to no
meetings. Afterwards $35$ of the meeting group and $20$ of the other group
passed all their classes.

\vspace{2pt}
(a) Find each group's pass percent and the gap.

\begin{work}
  \dfrac{35}{50} &= 0.70 = 70\% \\
  \dfrac{20}{50} &= 0.40 = 40\% \\
  70\% - 40\% &= 30 \text{ percentage points}
\end{work}

(b) Treatment: \ans{Required club meetings} \quad Control group: \ans{The 50 with no meetings}

\vspace{3pt}
(c) Instead, Oakmont \textbf{blocks} on last quarter's grades: $60$ volunteers
passed everything and $40$ did not. Randomly split each block in half. How many
go to each group?

\begin{work}
  60 \div 2 &= 30 \text{ per group (passed-everything block)} \\
  40 \div 2 &= 20 \text{ per group (did-not block)}
\end{work}

% ---------------------------------------------------------------- C8 (1.8)
\boxguard[30]
\item Westfield's student government asked a random sample of $50$ students how
to spend a \$$2{,}000$ grant. Complete the table, then complete the bar graph.

\vspace{2pt}
\renewcommand{\arraystretch}{1.25}
\begin{tabularx}{\linewidth}{@{} X c c @{}}
\rowcolor{frost}
\textbf{Choice} & \textbf{Count} & \textbf{Relative frequency} \\
\midrule
Spirit week     & $20$ & \ans{$40\%$} \\
Field trips     & $15$ & \ans{$30\%$} \\
New computers   & $10$ & \ans{$20\%$} \\
Longer lunch    & $5$  & \ans{$10\%$} \\
\end{tabularx}

\vspace{6pt}
\begin{center}
\begin{tikzpicture}[x=1cm, y=0.075cm]
  \foreach \y in {5,15,25,35}
    \draw[linegray, very thin] (0,\y) -- (10.4,\y);
  \foreach \y in {10,20,30,40}
    \draw[linegray] (0,\y) -- (10.4,\y);
  \foreach \y in {0,10,20,30,40}
    \node[left, font=\scriptsize] at (-0.05,\y) {\y\%};
  \foreach \x in {2.6,5.2,7.8}
    \draw[linegray, dashed] (\x,0) -- (\x,45);
  \foreach \i/\lab in {1/Spirit week, 2/Field trips, 3/New computers,
                       4/Longer lunch}
    \node[below, font=\scriptsize, align=center, text width=2.4cm]
      at ({2.6*\i-1.3},-0.5) {\lab};
  \draw[cerulean, thick] (0,45) -- (0,0) -- (10.4,0);
  % Key only: the four correct bars, drawn inside the existing bounding box.
  \foreach \i/\h in {1/40, 2/30, 3/20, 4/10}
    \fill[keyred, opacity=0.55] ({2.6*\i-2.0},0) rectangle ({2.6*\i-0.6},\h);
\end{tikzpicture}
\end{center}

\vspace{2pt}
About how many of all $800$ Westfield students would choose spirit week?

\begin{work}
  0.40 \times 800 &= 320 \text{ students}
\end{work}

\end{enumerate}
\normalsize

% ======================================================================
\boxguard[10]
\parthead{Part D --- Extended Response \hfill 3 questions, 10 points each}

\small
\begin{enumerate}[leftmargin=2.2em, itemsep=10pt, topsep=2pt, start=31]

% ---------------------------------------------------------------- D1
\item Oakmont has now run \textbf{two} studies of the same Homework Club.

\vspace{1pt}
\textbf{Study A.} A random sample of students was watched. Of the students who
\emph{chose} to join the club, $55\%$ passed all classes; of those who did not
join, $70\%$ did. \\
\textbf{Study B.} Of $100$ volunteers, half were \emph{assigned} to required
meetings. Of those, $70\%$ passed all classes; of the others, $40\%$ did.

\vspace{3pt}
(a) Which study is observational and which is an experiment? Give the one
feature that decides it.

\ansline{A is observational, B is an experiment. The deciding feature is who decides
the groups:}
\ansline{in A the students joined on their own; in B a random generator assigned them.}

(b) Study A's gap is $15$ points and Study B's is $30$ points. Which study lets
the school write \emph{``the Homework Club caused more students to pass''}?
Justify your choice.

\ansline{Study B. Random assignment makes the two groups alike at the start, so the
meetings are}
\ansline{the only difference left. Study A earns only ``is associated with,'' whatever
its gap.}

(c) Explain how the same Homework Club can look \emph{harmful} in Study A and
\emph{helpful} in Study B.

\ansline{In A the students who joined were already struggling, so their group started
behind.}
\ansline{In B nobody chose, so that head start cannot be hiding inside the gap.}

% ---------------------------------------------------------------- D2
\boxguard[26]
\item \textbf{Pinehurst Apartments} mailed a survey about a proposed community
garden to all $800$ households. $200$ were returned, and $128$ of those said yes.
The management office's own sign-up records show that $320$ of the $800$
households actually want one.

\vspace{2pt}
(a) Find the response rate, the survey's percent saying yes, and the number of
households that percent points to.

\begin{work}
  \dfrac{200}{800} &= 0.25 = 25\% \text{ response rate} \\
  \dfrac{128}{200} &= 0.64 = 64\% \\
  0.64 \times 800 &= 512 \text{ households}
\end{work}

(b) The records say the truth is $320$ of $800$, or $40\%$. Name the bias, and
give its size in percentage points \emph{and} in households.

\ansline{Nonresponse bias --- the $75\%$ who threw the survey away cared less about a
garden.}
\ansline{The survey overestimates by $24$ percentage points, or $512 - 320 = 192$
households.}

(c) The manager mails again and gets $400$ returns, $256$ of them yes. Compute
that percent and explain what the result shows.

\begin{work}
  \dfrac{256}{400} &= 0.64 = 64\%
\end{work}

\ansline{Twice the responses, the same wrong $64\%$: sample size does not remove bias.}

(d) Describe one change that would actually reduce the bias, and name the bias
it removes.

\ansline{Draw a random sample of households and follow up in person until nearly all
of them}
\ansline{answer. That removes the nonresponse bias, which a bigger mailing cannot.}

% ---------------------------------------------------------------- D3
\boxguard[24]
\item \textbf{Lakeside Farmers Market} wants to know whether to add a Wednesday
market. Its loyalty program has a numbered list of all $800$ regular shoppers,
and the three shopper types --- produce only, crafts only, and both --- are
believed to differ a lot in their answers.

\vspace{2pt}
(a) Choose \textbf{one} of the four sampling techniques and justify it from what
this context allows.

\ansline{Stratified. A full numbered list exists, and the three shopper types are
expected to}
\ansline{differ, so sampling within each type guarantees all three are represented in
proportion.}

(b) Name one of the five \textbf{data collection methods} and say why it fits
this question.

\ansline{Survey. The question has a small set of fixed answers and needs many
shoppers,}
\ansline{so a short written instrument can be counted quickly.}

(c) Instead, one vendor asks her own $20$ regular customers; $8$ say yes. She
reports: \emph{``$8 \div 20 = 40\%$, and $0.40 \times 800 = 320$ shoppers want
the Wednesday market.''} Her arithmetic is correct. Explain what is wrong with
the report, and state what she \emph{is} allowed to say.

\ansline{Her own customers are a convenience sample --- chance never chose them, and
the other}
\ansline{$780$ shoppers had no possibility of being asked, so the $320$ is not an
estimate for the market.}
\ansline{She may say only that $40\%$ of her own regular customers want the Wednesday
market.}

\end{enumerate}
\normalsize

\end{document}
